\chapter{表格统计学习答案}
\begin{enumerate}[leftmargin=*]
  \item 树无需标准化、能表示阈值和交互、原生支持多种缺失策略，适合中等样本表格；相对线性模型失去简单全局系数，相对神经网络不擅长端到端学习序列、文本和图像表示。
  \item CART 逐节点最小化不纯度；随机森林平均 bootstrap 树并随机抽特征以降相关；gradient boosting 在函数空间拟合负梯度；XGBoost 再加入二阶曲率、叶权重和结构惩罚。
  \item $s=-1.5,0,2$ 的 SSE 分别为 $17.33,10,10.67$（按各侧均值计算），所以 $s=0$ 最优；不分裂 SSE 为 46。
  \item 叶目标为 $G_jw+\frac12(H_j+\lambda)w^2+\gamma$。一阶条件 $G_j+(H_j+\lambda)w=0$ 给出 $w_j^*=-G_j/(H_j+\lambda)$。
  \item 使用相同外层 walk-forward，内层选择复杂度；报告模型、时间、验证损失、事先确定仓位映射后的净收益与换手。若 MLP 不能稳定胜过提升树，应保留提升树。
  \item 随机拆行会让同一资产相邻日期、同一市场冲击和重叠标签同时进入训练测试；AUC 衡量的是污染后的区分能力。应按决策时间切分并 purge 标签区间。
\end{enumerate}

\section*{逐题推导与核对}
\begin{enumerate}[leftmargin=*]
  \item 树模型按阈值递归切分，不需要先把特征缩放到同一量纲，并能直接表示非线性和交互；代价是全局系数不如线性模型透明，外推和高维稀疏表示也较弱。神经网络能端到端学习表示，但需要更多样本、调参和正则化；因此树是量化表格数据的稳健基线而不是自动最优答案。
  \item CART 在每个候选节点选择使子节点不纯度下降最大的切分；随机森林通过 bootstrap 和特征子采样降低树间相关，再平均；boosting 逐轮拟合当前残差/负梯度；XGBoost 用二阶 Taylor 展开和叶结构惩罚。比较时要把随机化对象、损失近似和停止规则分别写清。
  \item 候选阈值 $s=-1.5,0,2$ 的左右均值分别代入 $\sum_{i\in L}(y_i-\bar y_L)^2+\sum_{i\in R}(y_i-\bar y_R)^2$，得到 $17.33,10,10.67$；不切分时总 SSE 为 $46$。因此 $s=0$ 最优，但仍需检查叶节点最小样本数和 tie-breaking 规则。
  \item 对叶 $j$ 的二阶近似为 $G_jw+\frac12(H_j+\lambda)w^2+\gamma$。求导得到 $G_j+(H_j+\lambda)w=0$，因 $H_j+\lambda>0$，唯一极小点为 $w_j^*=-G_j/(H_j+\lambda)$；代回目标得该叶的增益，再与分裂前目标比较是否值得切分。
  \item 在同一 walk-forward 外层切分内，训练集再做内层选择；四个模型均固定特征、目标、随机种子和仓位映射。表格至少列训练时间、验证 MSE、净收益、换手和区间；若 MLP 只在一次窗口胜出，不能据此推翻更稳定的提升树。
  \item 随机拆行会使同一资产的相邻日期、同一冲击和重叠标签两边同时出现。AUC 因此估计的是污染后的可区分性，而不是未来预测能力；修复为按时间切分、按资产/事件 purge，并把 scaler 和特征选择限制在训练窗内。
\end{enumerate}
